\documentclass[11pt,letterpaper]{article}
\usepackage[spanish,es-nodecimaldot]{babel}
\usepackage[utf8]{inputenc}
\usepackage[T1]{fontenc}
\usepackage{lmodern}
\usepackage{geometry}
\geometry{left=2.5cm,right=2.5cm,top=2.4cm,bottom=2.4cm}
\usepackage{amsmath,amssymb,mathtools,bm}
\usepackage{booktabs,array,tabularx,multirow}
\usepackage{graphicx,float,caption}
\usepackage{microtype}
\usepackage{siunitx}
\usepackage{xcolor}
\usepackage{hyperref}
\usepackage{fancyhdr}
\usepackage{enumitem}
\usepackage{listings}
\usepackage{tikz}
\usetikzlibrary{arrows.meta,positioning}
\hypersetup{colorlinks=true,linkcolor=blue!55!black,urlcolor=blue!55!black,citecolor=blue!55!black}
\sisetup{output-decimal-marker={,},group-separator={\,}}
\setlength{\parindent}{0pt}
\setlength{\parskip}{0.45em}
\pagestyle{fancy}
\setlength{\headheight}{14pt}
\fancyhf{}
\lhead{Control adaptativo MRAC}
\rhead{JAXA Mu-pal}
\cfoot{\thepage}
\renewcommand{\headrulewidth}{0.3pt}
\newcommand{\norm}[1]{\left\lVert #1\right\rVert}
\newcommand{\figuraopcional}[3]{%
\begin{figure}[htbp]
    \centering
    \IfFileExists{#1}{\includegraphics[width=0.92\textwidth]{#1}}{%
        \fbox{\begin{minipage}[c][3.6cm][c]{0.86\textwidth}\centering
            \textbf{Espacio reservado para la gráfica de MATLAB}\\[0.5em]
            Archivo: \texttt{\detokenize{#1}}\\[0.5em]
            Exportar la figura desde MATLAB y cargarla en Overleaf.
        \end{minipage}}}
    \caption{#2}
    \label{#3}
\end{figure}}
\lstdefinestyle{matlab}{language=Matlab,basicstyle=\ttfamily\footnotesize,keywordstyle=\color{blue!60!black},commentstyle=\color{green!40!black},stringstyle=\color{orange!50!black},frame=single,breaklines=true,showstringspaces=false,columns=fullflexible}

\begin{document}
\begin{titlepage}
    \centering
    {\Large Universidad Nacional de Colombia\par}
    \vspace{1.2cm}
    {\large Informe de taller\par}
    \vspace{0.5cm}
    {\LARGE\bfseries Control adaptativo con modelo de referencia (MRAC directo)\par}
    \vspace{0.5cm}
    {\Large Aeronave JAXA Mu-pal: dinámica longitudinal\par}
    \vspace{1.8cm}
    \textbf{Estudiante(s):} \underline{\hspace{8cm}}\par
    \vspace{0.6cm}
    \textbf{Asignatura:} \underline{\hspace{8cm}}\par
    \vspace{0.6cm}
    \textbf{Docente:} \underline{\hspace{8cm}}\par
    \vfill
    Bogotá, Colombia\\
    Noviembre de 2026
\end{titlepage}

\begin{abstract}
Se estudia un controlador adaptativo directo con modelo de referencia (MRAC) para la dinámica longitudinal de la aeronave JAXA Mu-pal, con el objetivo de regular el ángulo de cabeceo y mantener una evolución estable de los estados. Se comprueba la estabilidad del modelo de referencia, se obtienen ganancias de acoplamiento aproximadas mediante pseudoinversa y se derivan las leyes adaptativas utilizando una función de Lyapunov. Se analizan dos ensayos numéricos implementados en MATLAB, uno con entrada escalón y otro con tren de pulsos. Los residuos de acoplamiento son pequeños, pero diferentes de cero; además, las ganancias adaptadas al final de las simulaciones no reproducen la matriz de referencia. Por ello se distinguen las garantías teóricas, válidas bajo acoplamiento exacto, de las observaciones obtenidas para el problema numérico concreto.
\end{abstract}

\section{Objetivos}
\textbf{Objetivo general.} Diseñar y evaluar un controlador MRAC directo para el sistema longitudinal JAXA Mu-pal a partir de un modelo de referencia establecido.

\textbf{Objetivos específicos.}
\begin{enumerate}[itemsep=2pt]
    \item Analizar la estabilidad del modelo de referencia mediante los autovalores de $A_m$.
    \item Calcular las ganancias ideales aproximadas y evaluar las condiciones de acoplamiento.
    \item Derivar las leyes de actualización adaptativa con el método directo de Lyapunov.
    \item Comparar el seguimiento y las ganancias para entradas escalón y tren de pulsos.
    \item Interpretar las diferencias entre $A+B\hat K_x^T$ y $A_m$.
\end{enumerate}

\section{Modelo matemático de la aeronave}
De acuerdo con el enunciado del taller \cite{sofrony2021}, la dinámica longitudinal linealizada se representa como
\begin{equation}\label{eq:planta}
    \dot x(t)=Ax(t)+Bu(t), \qquad
    x=\begin{bmatrix}u_x&u_z&\theta&q\end{bmatrix}^{T},
\end{equation}
donde $u_x$ y $u_z$ son componentes de velocidad, $\theta$ es el ángulo de cabeceo (\emph{pitch}), $q$ su velocidad angular y $u(t)$ la señal de control.
\begin{equation}\label{eq:AB}
A=\begin{bmatrix}
-0.0175&0.173&-9.77&-5.63\\
-0.192&-1.09&-0.846&64.6\\
0&0&0&1\\
0.0081&-0.0738&0.0062&-1.9
\end{bmatrix},\qquad
B=\begin{bmatrix}-0.428\\4.91\\0\\4.22\end{bmatrix}.
\end{equation}
Se desea que la planta siga un modelo estable descrito por
\begin{equation}\label{eq:referencia}
    \dot x_m(t)=A_m x_m(t)+B_m r(t),
\end{equation}
con matrices prescritas en el taller:
\begin{equation}\label{eq:AmBm}
A_m=\begin{bmatrix}
-0.03&0.166&125.6&37.29\\
-0.052&-1.02&-1554.7&-427.82\\
0&0&0&1\\
0.128&-0.0142&-1335.49&-425.12
\end{bmatrix},\quad
B_m=\begin{bmatrix}-138.1\\1584.2\\0\\1361.6\end{bmatrix}.
\end{equation}
\textbf{Convención de error:}
\begin{equation}\label{eq:error}
    e(t)=x_m(t)-x(t).
\end{equation}
El objetivo de seguimiento es que $e(t)$ tienda a cero, particularmente su componente asociada al ángulo $\theta$.

\section{Paso 1. Modelo de referencia y ganancias de acoplamiento}
\subsection{Estabilidad del modelo de referencia}
La solución homogénea del modelo de referencia viene dada por $x_m(t)=e^{A_mt}x_m(0)$. Para que el término asociado a las condiciones iniciales decaiga, $A_m$ debe ser una matriz Hurwitz, es decir, todos sus autovalores deben tener parte real negativa.

Los autovalores calculados en MATLAB fueron
\begin{equation}\label{eq:eigAm}
    \lambda(A_m)=\{-421.9808,\ -3.1703,\ -0.0528,\ -0.9661\}.
\end{equation}
Dado que todos son negativos, el modelo de referencia es asintóticamente estable. La presencia de polos con magnitudes muy diferentes (en particular, $-421.9808$ y $-0.0528$) evidencia escalas temporales muy distintas, lo cual puede hacer rígida la integración numérica.

En respuesta a una referencia constante $r(t)=r_0$ y si $x_m(0)=0$, la respuesta se expresa como
\begin{equation}
    x_m(t)=\int_0^t e^{A_m(t-\tau)}B_mr_0\,\mathrm d\tau,
    \qquad x_m(\infty)=-A_m^{-1}B_mr_0.
\end{equation}
Este resultado permite interpretar la simulación ante un escalón sin confundir la respuesta del modelo con la de la planta adaptativa.

\subsection{Condiciones de acoplamiento}
Se adopta la ley de control de estructura fija y parámetros ajustables:
\begin{equation}\label{eq:leycontrol}
    u(t)=\hat K_x^T(t)x(t)+\hat k_r(t)r(t),
\end{equation}
donde $\hat K_x\in\mathbb R^4$ y $\hat k_r\in\mathbb R$. Sustituyendo \eqref{eq:leycontrol} en \eqref{eq:planta}:
\begin{align}
    \dot x &= Ax+B\big(\hat K_x^T x+\hat k_r r\big)\\
           &=\big(A+B\hat K_x^T\big)x+B\hat k_r r.
\end{align}
Para reproducir exactamente el modelo de referencia con ganancias constantes $K_x^*$ y $k_r^*$, deben satisfacerse simultáneamente
\begin{equation}\label{eq:matching}
    \boxed{A+BK_x^{*T}=A_m,\qquad Bk_r^*=B_m.}
\end{equation}
Equivalentemente,
\begin{equation}
    BK_x^{*T}=A_m-A,\qquad Bk_r^*=B_m.
\end{equation}

\subsection{Cálculo de ganancias aproximadas mediante pseudoinversa}
Como $B$ es un vector de cuatro filas y una columna, las cuatro columnas de $A_m-A$ deben pertenecer al subespacio generado por $B$ para cumplir exactamente la primera condición. Esto no está garantizado por los datos impresos.

Para calcular el mejor ajuste en mínimos cuadrados se usa la pseudoinversa de Moore--Penrose:
\begin{equation}
 B^{\dagger}=(B^TB)^{-1}B^T,\qquad
 K_x^{*T}=B^{\dagger}(A_m-A),\qquad k_r^*=B^{\dagger}B_m.
\end{equation}
El denominador escalar es
\begin{equation}
    B^TB=(-0.428)^2+(4.91)^2+0^2+(4.22)^2=42.099684.
\end{equation}
Por ejemplo, para la tercera componente de $K_x^*$, correspondiente a $\theta$, se toma la tercera columna de $A_m-A$:
\begin{equation}
    \Delta a_3=
    \begin{bmatrix}135.37&-1553.854&0&-1335.4962\end{bmatrix}^{T}
\end{equation}
y se obtiene
\begin{align}
    k_{x,3}^*&=\frac{B^T\Delta a_3}{B^TB}\\
    &=\frac{(-0.428)(135.37)+(4.91)(-1553.854)+(4.22)(-1335.4962)}{42.099684}\\
    &\approx -316.4669.
\end{align}
Repitiendo el procedimiento para las demás columnas y para $B_m$ se encuentran los valores reportados por MATLAB:
\begin{equation}\label{eq:Kstar}
    \boxed{K_x^{*T}\approx
    \begin{bmatrix}0.0285&0.0142&-316.4669&-100.2891\end{bmatrix},}
\end{equation}
\begin{equation}\label{eq:krstar}
    \boxed{k_r^*\approx 322.65042179.}
\end{equation}

\subsection{Verificación de las condiciones de acoplamiento}
Se definen los residuos de acoplamiento
\begin{equation}\label{eq:residuos}
    \Delta_A=A_m-(A+BK_x^{*T}),\qquad \Delta_B=B_m-Bk_r^*.
\end{equation}
Los resultados fueron
\begin{equation}\label{eq:residuosnumeric}
    \boxed{\norm{\Delta_A}_F=0.078168555,\qquad
    \norm{\Delta_B}_2=0.021151833.}
\end{equation}
Ambas cantidades son diferentes de cero: \textbf{las matrices dadas permiten un acoplamiento numérico muy cercano, pero no exacto}. Una explicación plausible es el redondeo de las entradas de $A_m$ y $B_m$; sin matrices de mayor precisión no es posible confirmar que esta sea la única causa. La demostración teórica estándar debe interpretarse bajo la hipótesis de acoplamiento exacto.

\section{Paso 2. Diseño de leyes adaptativas mediante Lyapunov}
\subsection{Dinámica del error en el caso ideal}
Definamos los errores paramétricos con respecto a las ganancias ideales:
\begin{equation}\label{eq:paramerror}
    \tilde K_x=\hat K_x-K_x^*,\qquad
    \tilde k_r=\hat k_r-k_r^*.
\end{equation}
Al derivar \eqref{eq:error} y sustituir las dinámicas del modelo y la planta:
\begin{align}
    \dot e&=\dot x_m-\dot x\\
    &=A_m x_m+B_mr-Ax-B\big(\hat K_x^Tx+\hat k_rr\big).
\end{align}
Usando $x_m=x+e$ y las igualdades exactas \eqref{eq:matching}:
\begin{align}
    \dot e&=A_me+(A_m-A)x+(B_m-B\hat k_r)r-B\hat K_x^Tx\\
    &=A_me+B(K_x^*-\hat K_x)^Tx+B(k_r^*-\hat k_r)r.
\end{align}
Así,
\begin{equation}\label{eq:errorideal}
    \boxed{\dot e=A_me-B\big(\tilde K_x^Tx+\tilde k_r r\big).}
\end{equation}

\subsection{Matriz de Lyapunov}
Como $A_m$ es Hurwitz, para toda matriz simétrica definida positiva $Q=Q^T\succ0$ existe una única matriz simétrica definida positiva $P=P^T\succ0$ que satisface
\begin{equation}\label{eq:lyapeq}
    \boxed{A_m^TP+PA_m=-Q.}
\end{equation}
En MATLAB, para un valor seleccionado de $Q$, se resuelve con \texttt{P = lyap(Am',Q)}. Los autovalores de $P$ informados por MATLAB (visualizados con factor común $10^5$) son
\begin{equation}\label{eq:eigP}
    \lambda(P)=10^5\times
    \begin{bmatrix}0.0000&0.0002&0.0083&2.3630\end{bmatrix}^{T},
\end{equation}
donde los ceros son consecuencia del formato de redondeo de MATLAB, no evidencia de autovalores exactamente nulos. El cálculo numérico proporciona $\norm{A_m^TP+PA_m+Q}\approx 3.388\times10^{-6}$. El valor bajo de este residuo apoya que la ecuación de Lyapunov se resolvió adecuadamente, a la precisión de la simulación. La elección exacta de $Q$ debe conservarse de acuerdo con el script ejecutado.

\subsection{Construcción de la función candidata}
Se consideran parámetros de adaptación $\Gamma_x=\Gamma_x^T\succ0$ y $\gamma_r>0$, y la función candidata
\begin{equation}\label{eq:V}
    \boxed{V=e^TPe+\tilde K_x^T\Gamma_x^{-1}\tilde K_x
          +\frac{\tilde k_r^2}{\gamma_r}.}
\end{equation}
Dado que $P$ y $\Gamma_x$ son definidas positivas, y $\gamma_r>0$, se cumple $V\geq0$ y $V=0$ únicamente si $e=0$, $\tilde K_x=0$ y $\tilde k_r=0$.

Derivando respecto del tiempo, suponiendo constantes las ganancias ideales,
\begin{align}
    \dot V&=\dot e^TPe+e^TP\dot e
    +2\tilde K_x^T\Gamma_x^{-1}\dot{\hat K}_x
    +\frac{2\tilde k_r}{\gamma_r}\dot{\hat k}_r\\
    &=2e^TP\dot e+2\tilde K_x^T\Gamma_x^{-1}\dot{\hat K}_x
    +\frac{2\tilde k_r}{\gamma_r}\dot{\hat k}_r.
\end{align}
Sustituyendo \eqref{eq:errorideal}:
\begin{align}
    \dot V={}&e^T(A_m^TP+PA_m)e
    -2e^TPB\,\tilde K_x^Tx-2e^TPB\,\tilde k_r r\\
    &+2\tilde K_x^T\Gamma_x^{-1}\dot{\hat K}_x
    +\frac{2\tilde k_r}{\gamma_r}\dot{\hat k}_r.
\end{align}
Aplicando \eqref{eq:lyapeq}:
\begin{align}\label{eq:Vdotexpanded}
    \dot V={}&-e^TQe
    +2\tilde K_x^T\Big(\Gamma_x^{-1}\dot{\hat K}_x-xe^TPB\Big)\\
    &+2\tilde k_r\Big(\frac{\dot{\hat k}_r}{\gamma_r}-re^TPB\Big).
\end{align}
Los términos cruzados se cancelan con las siguientes leyes:
\begin{equation}\label{eq:adapt}
    \boxed{\dot{\hat K}_x=\Gamma_x x\,(e^TPB),\qquad
    \dot{\hat k}_r=\gamma_r r\,(e^TPB).}
\end{equation}
Nótese que los signos son \textbf{positivos} porque se ha definido $e=x_m-x$.
Sustituyendo \eqref{eq:adapt} en \eqref{eq:Vdotexpanded},
\begin{equation}\label{eq:Vdot}
    \boxed{\dot V=-e^TQe\leq0.}
\end{equation}

\subsection{Alcance de la demostración}
Para las \emph{condiciones de acoplamiento exactas} y una referencia acotada, la desigualdad \eqref{eq:Vdot} garantiza acotamiento del error y de los errores paramétricos. Además, usando la estabilidad de $A_m$, las señales permanecen acotadas y, bajo las condiciones habituales de regularidad y acotamiento de $\dot e$, puede aplicarse el lema de Barbalat para concluir $e(t)\to0$ cuando $t\to\infty$. Esta conclusión no implica necesariamente que las ganancias estimadas converjan a las ganancias ideales: la convergencia paramétrica puede requerir excitación persistente.

\textbf{Limitación para los datos del taller.} Como $\Delta_A\neq0$ y $\Delta_B\neq0$, la dinámica real del error contiene términos adicionales:
\begin{equation}\label{eq:realerror}
    \boxed{\dot e=A_me-B(\tilde K_x^Tx+\tilde k_rr)+\Delta_Ax+\Delta_Br.}
\end{equation}
Con las mismas leyes adaptativas se obtiene
\begin{equation}\label{eq:realVdot}
    \boxed{\dot V=-e^TQe+2e^TP(\Delta_Ax+\Delta_Br).}
\end{equation}
El segundo término puede tener cualquier signo. Por tanto, \emph{con los datos redondeados del taller la demostración anterior no garantiza por sí sola que $e(t)\to0$}, aunque sí describe correctamente el caso de acoplamiento ideal.

\section{Implementación numérica en MATLAB}
\subsection{Estados aumentados y condiciones iniciales}
La simulación implementa simultáneamente el modelo de la planta, el modelo de referencia y las leyes de adaptación, integrando un sistema aumentado de 13 variables:
\begin{equation}
    z=\begin{bmatrix}x^T & x_m^T &\hat K_x^T &\hat k_r\end{bmatrix}^{T}\in\mathbb R^{13}.
\end{equation}
La asignación utilizada es
\begin{equation}
\begin{aligned}
    x&=z_{1:4}, & x_m&=z_{5:8},\\
    \hat K_x&=z_{9:12}, & \hat k_r&=z_{13}.
\end{aligned}
\end{equation}
Se utilizó el vector inicial nulo $z(0)=0_{13\times1}$ y el integrador \texttt{ode15s}, apropiado para explorar sistemas con escalas temporales muy diferentes. Las simulaciones se ejecutaron por separado para una referencia escalón y otra de tren de pulsos. Como la impresión numérica proporcionada no incluye el tiempo final ni los parámetros temporales y amplitudes de estas señales, dichos valores deben consultarse directamente en el script antes de describirlos cuantitativamente en el texto o las figuras.

La ganancia de adaptación vectorial y escalar se parametriza mediante $\Gamma_x$ y $\gamma_r$; se recomienda reportar aquí los valores exactos que quedaron definidos en la versión de MATLAB ejecutada. La estructura central de la función de dinámica es:
\begin{lstlisting}[style=matlab]
x    = z(1:4);        xm   = z(5:8);
Khat = z(9:12);       krhat = z(13);
e = xm - x;
u = Khat.'*x + krhat*r;
dx = A*x + B*u;
dxm = Am*xm + Bm*r;
dKhat = GammaX*x*(e.'*P*B);
dkrhat = gammaR*r*(e.'*P*B);
dz = [dx; dxm; dKhat; dkrhat];
\end{lstlisting}

\section{Resultados y discusión}
\subsection{Ensayo 1: entrada escalón}
Los valores finales informados fueron
\begin{equation}\label{eq:stepK}
\hat K_x^T(t_f)=\begin{bmatrix}-0.8660&-9.0106&3.9486&-239.4407\end{bmatrix},
\qquad \hat k_r(t_f)=168.36175915.
\end{equation}
La norma del error de estados en el instante final es
\begin{equation}
    \norm{e(t_f)}_2=0.19759963.
\end{equation}
Por consiguiente, el error al final de la simulación \emph{no} es nulo. Como $\norm{e}_2$ combina estados con magnitudes y unidades distintas, resulta conveniente complementar esta métrica con la gráfica de $e_\theta(t)=\theta_m(t)-\theta(t)$ para evaluar específicamente el control del cabeceo.

La matriz instantánea de lazo cerrado evaluada con las ganancias finales fue
\begin{equation}\label{eq:aclstep}
    A+B\hat K_x^T(t_f)\approx10^3\begin{bmatrix}
     0.0004& 0.0040&-0.0115& 0.0969\\
    -0.0044&-0.0453& 0.0185&-1.1111\\
     0&0&0&0.0010\\
    -0.0036&-0.0381& 0.0167&-1.0123
    \end{bmatrix}.
\end{equation}
Comparándola con $A_m$ se obtuvo
\begin{equation}\label{eq:stepResid}
    \boxed{\norm{A_m-(A+B\hat K_x^T(t_f))}_F=2267.3457.}
\end{equation}
La diferencia es muy superior al residuo obtenido con $K_x^*$ en la sección~\ref{eq:residuosnumeric}. Esto muestra que la simulación no terminó en un estado de acoplamiento cercano al ideal, aunque la norma del error de estados final resultó comparativamente moderada. No debe confundirse seguimiento transitorio con identificación exacta de ganancias.

\figuraopcional{figuras/estados_escalon.png}{Comparación temporal entre estados de la planta $x$ y del modelo de referencia $x_m$ para el ensayo de escalón.}{fig:estados_escalon}
\figuraopcional{figuras/ganancias_escalon.png}{Evolución temporal de $\hat K_x(t)$ y $\hat k_r(t)$ bajo referencia escalón.}{fig:ganancias_escalon}

\subsection{Ensayo 2: tren de pulsos}
Los valores finales reportados fueron
\begin{equation}\label{eq:pulseK}
\hat K_x^T(t_f)=\begin{bmatrix}-10.6056&-51.9365&1.3499&-235.3359\end{bmatrix},
\qquad \hat k_r(t_f)=168.32289102.
\end{equation}
La norma del error final fue
\begin{equation}
    \norm{e(t_f)}_2=0.20744758.
\end{equation}
La matriz instantánea de lazo cerrado, con las ganancias finales congeladas, fue
\begin{equation}\label{eq:aclpulse}
    A+B\hat K_x^T(t_f)\approx10^3\begin{bmatrix}
     0.0045&0.0224&-0.0103&0.0951\\
    -0.0523&-0.2561&0.0058&-1.0909\\
     0&0&0&0.0010\\
    -0.0447&-0.2192&0.0057&-0.9950
    \end{bmatrix}.
\end{equation}
La discrepancia respecto a la matriz de referencia fue
\begin{equation}\label{eq:pulseResid}
    \boxed{\norm{A_m-(A+B\hat K_x^T(t_f))}_F=2266.8403.}
\end{equation}
Este valor es ligeramente menor que el obtenido con el escalón, pero continúa siendo alto. La mayor variación observada en algunas componentes de $\hat K_x$ no implica, por sí sola, una mejora suficiente del acoplamiento ni de la convergencia del error.

\figuraopcional{figuras/estados_pulsos.png}{Comparación temporal entre los estados de planta y referencia ante un tren de pulsos.}{fig:estados_pulsos}
\figuraopcional{figuras/ganancias_pulsos.png}{Evolución temporal de las ganancias adaptativas ante un tren de pulsos.}{fig:ganancias_pulsos}

\subsection{Comparación cuantitativa de resultados}
\begin{table}[H]
\centering
\caption{Comparación de las simulaciones de MRAC directo.}
\label{tab:comparacion}
\small
\begin{tabularx}{\textwidth}{@{}Xrr@{}}
\toprule
\textbf{Magnitud} & \textbf{Escalón} & \textbf{Tren de pulsos}\\
\midrule
$\hat k_r(t_f)$ & $168.36175915$ & $168.32289102$\\
$\norm{e(t_f)}_2$ & $0.19759963$ & $0.20744758$\\
$\norm{A_m-(A+B\hat K_x^T(t_f))}_F$ & $2267.3457$ & $2266.8403$\\
\bottomrule
\end{tabularx}
\end{table}
En comparación, las ganancias de ajuste ideales aproximadas dan $k_r^*\approx322.6504$ y $\norm{\Delta_A}_F\approx0.07817$. Las normas del error final de ambas simulaciones son del mismo orden; el ensayo con escalón presenta un error final ligeramente menor, mientras que el tren de pulsos presenta una discrepancia de matrices ligeramente menor. Estas observaciones corresponden a un instante final y no prueban por sí mismas convergencia asintótica.

\subsection{Análisis adicional de polos de las matrices finales}
Como comprobación adicional, es posible calcular los autovalores de $A+B\hat K_x^T(t_f)$ utilizando las ganancias finales \emph{redondeadas} que se imprimieron en MATLAB. Esta operación equivale a congelar momentáneamente las ganancias, y no debe confundirse con un análisis completo de estabilidad del sistema adaptativo, cuyas ganancias cambian con el tiempo.
\begin{table}[H]
\centering
\caption{Autovalores aproximados de las matrices con ganancias finales congeladas.}
\label{tab:autovalores_finales}
\small
\begin{tabular}{@{}ll@{}}
\toprule
\textbf{Ensayo} & \textbf{Autovalores aproximados}\\
\midrule
Escalón & $-1053.9573,\ -3.3442,\ -0.00858\pm0.10252\,\mathrm{i}$\\
Tren de pulsos & $-1234.1071,\ -12.4768,\ -0.03892,\ +0.02880$\\
\bottomrule
\end{tabular}
\end{table}
En el caso del escalón los polos de la matriz congelada tienen parte real negativa, aunque aparece un par complejo débilmente amortiguado. En el caso del tren de pulsos aparece un polo real positivo con las ganancias finales redondeadas. Esto advierte que \emph{mantener fijas} estas ganancias del segundo ensayo no produciría una matriz Hurwitz. Sin embargo, no basta para afirmar que toda la trayectoria adaptativa sea inestable, porque $\hat K_x(t)$ evoluciona durante la simulación.

\subsection{Interpretación de las diferencias entre los ensayos}
\begin{enumerate}[leftmargin=*,itemsep=4pt]
\item \textbf{Seguimiento frente a convergencia paramétrica.} Un error relativamente pequeño en un instante no exige que $\hat K_x$ y $\hat k_r$ hayan convergido a sus valores ideales. La adaptación depende de la información presente en $x(t)$, $r(t)$ y $e(t)$.
\item \textbf{Excitación de la referencia.} El tren de pulsos modifica repetidamente la entrada y puede excitar componentes diferentes del sistema, pero no es suficiente afirmar que aporta excitación persistente sin verificar sus propiedades y el intervalo temporal.
\item \textbf{Influencia de los residuos.} Los errores no nulos $\Delta_A$ y $\Delta_B$ impiden aplicar literalmente la prueba $\dot V=-e^TQe$ al modelo numérico tal como está definido.
\item \textbf{Horizonte de simulación y selección de ganancias.} Los resultados finales también dependen del horizonte, las condiciones iniciales y los parámetros $\Gamma_x$, $\gamma_r$ y $Q$. Sin los gráficos completos no es posible establecer si el error se encontraba decreciendo, oscilando o creciendo cerca de $t_f$.
\end{enumerate}

\section{Conclusiones}
\begin{enumerate}[leftmargin=*,itemsep=5pt]
\item El modelo de referencia es asintóticamente estable, porque sus cuatro autovalores tienen parte real negativa. Las escalas temporales muy distintas justifican prestar atención a la elección del integrador numérico.
\item Se calcularon ganancias ideales aproximadas $K_x^{*T}=[0.0285,\ 0.0142,\ -316.4669,\ -100.2891]$ y $k_r^*=322.6504$. Los residuos de acoplamiento $0.07817$ y $0.02115$ confirman que los datos no satisfacen las igualdades exactamente.
\item Con la convención $e=x_m-x$, las leyes $\dot{\hat K}_x=\Gamma_xx(e^TPB)$ y $\dot{\hat k}_r=\gamma_rr(e^TPB)$ cancelan los términos cruzados en la derivada de Lyapunov, dando $\dot V=-e^TQe$ \emph{solo bajo acoplamiento exacto}.
\item En las simulaciones las normas finales del error fueron $0.19760$ para escalón y $0.20745$ para pulsos. En ninguno de los dos ensayos el valor final informado fue cero.
\item Las discrepancias $\norm{A_m-A-B\hat K_x^T(t_f)}_F$ fueron aproximadamente $2267.35$ y $2266.84$, respectivamente. Por ello no puede concluirse que las ganancias adaptativas finales hayan reproducido la dinámica deseada.
\item Para evaluar el desempeño del controlador en el objetivo principal del taller conviene observar individualmente $\theta(t)$, $\theta_m(t)$ y $e_\theta(t)$, además de la estabilidad numérica y de la evolución temporal de todas las ganancias.
\end{enumerate}

\section*{Recomendaciones para completar las figuras}
En el proyecto de Overleaf se debe crear una carpeta \texttt{figuras/} y cargar allí cuatro imágenes con los nombres empleados en las llamadas de las secciones de resultados. Las imágenes deben salir de las simulaciones efectivamente ejecutadas en MATLAB; los recuadros del archivo actual son marcadores de posición, \emph{no} gráficas sintéticas ni resultados calculados nuevamente. Para exportar la figura activa desde MATLAB:
\begin{lstlisting}[style=matlab]
if ~exist('figuras','dir'), mkdir('figuras'); end
exportgraphics(gcf,'figuras/estados_escalon.png', ...
    'Resolution',300);
\end{lstlisting}
El nombre del archivo debe modificarse para las restantes figuras, después de activar cada figura correspondiente. Es aconsejable incluir también gráficos del error de cabeceo y del control $u(t)$ si se dispone de ellos.

\begin{thebibliography}{9}
\bibitem{sofrony2021}
Sofrony, J. (2021, 9 de julio).
\textit{Control adaptativo con modelo de referencia}.
Guía del taller suministrada para la asignatura.
\bibitem{astromwittenmark}
Åström, K. J. y Wittenmark, B. (1995).
\textit{Adaptive Control} (2.\textsuperscript{a} ed.). Addison-Wesley.
\bibitem{narendranannaswamy}
Narendra, K. S. y Annaswamy, A. M. (1989).
\textit{Stable Adaptive Systems}. Prentice Hall.
\end{thebibliography}
\end{document}
